\documentclass[12pt]{article}
% Préambule commun à tous les documents extraits de « La Revue CAPES »
\usepackage[T1]{fontenc}
\usepackage[utf8]{inputenc}
\usepackage{lmodern}
\usepackage{amsmath,amssymb,amsthm,amsfonts,mathrsfs}
\usepackage{setspace}
\usepackage{listings}
\usepackage{pgfplots}
\pgfplotsset{compat=1.18}
\usepackage{longtable}
\usepackage{adjustbox}
\usepackage{lettrine}
\usepackage{tkz-tab}
\usepackage{changepage}
\usepackage{pdflscape}
\usepackage{graphicx}
\usepackage{tikz}
\usepackage{xcolor}
\usepackage[a4paper,top=2cm,bottom=2.5cm,left=2.5cm,right=2cm]{geometry}
\usepackage{fancyhdr}
\usepackage{draftwatermark}
\SetWatermarkText{La RevueCapes}
\SetWatermarkLightness{0.9}
\SetWatermarkScale{2}
\usepackage{hyperref}
\hypersetup{colorlinks=true,linkcolor=blue!50!black,urlcolor=blue!50!black}

\definecolor{revue}{RGB}{20,60,120}

\pagestyle{fancy}
\fancyhf{}
\renewcommand{\headrulewidth}{0.4pt}
\setlength{\headheight}{15pt}

% \entete{Matière}{Titre}{Type de document}
\newcommand{\entete}[3]{%
  \fancyhead[L]{\small\textsc{La Revue CAPES} -- #1}%
  \fancyhead[R]{\small #2}%
  \fancyfoot[C]{\thepage}%
  \thispagestyle{plain}%
  \begin{center}
    {\color{revue}\rule{\textwidth}{1.2pt}}\\[4pt]
    {\small\textsc{Concours directs CAP-CEG / CAPES -- Burkina Faso}}\\[6pt]
    {\LARGE\bfseries\color{revue} #2}\\[6pt]
    {\large\itshape #1 -- #3}\\[2pt]
    {\color{revue}\rule{\textwidth}{1.2pt}}
  \end{center}
  \vspace{0.5cm}%
}

% Encadré pour les remarques ajoutées dans les corrigés
\newcommand{\remarque}[1]{\par\medskip\noindent\fcolorbox{revue}{revue!6}{\parbox{\dimexpr\linewidth-2\fboxsep-2\fboxrule}{\textbf{\color{revue}Remarque.} #1}}\par\medskip}


\begin{document}
\entete{Culture générale}{Corrigé du sujet type de dissertation n°9}{Correction}
\begin{spacing}{1.5}

\noindent\textbf{Sujet.} La démocratie est-elle réellement la forme de gouvernement la plus adaptée aux pays africains ?

\rubrique{1. Analyse du sujet}

\begin{itemize}
\item \textbf{Type de sujet :} question fermée, nuancée par l'adverbe « réellement », qui invite à interroger une idée reçue. Plan dialectique attendu.
\item \textbf{Mots-clés :}
  \begin{itemize}
  \item « démocratie » : régime fondé sur la souveraineté du peuple, les élections libres, le pluralisme, la séparation des pouvoirs et le respect des droits ;
  \item « forme de gouvernement la plus adaptée » : comparaison avec d'autres régimes (monarchie, régimes militaires, partis uniques) et adéquation aux réalités sociales, culturelles et économiques ;
  \item « pays africains » : diversité de situations, mais aussi des points communs (héritage colonial, jeunesse des États, pauvreté, insécurité dans certaines régions).
  \end{itemize}
\item \textbf{Problématique :} la démocratie, souvent présentée comme un modèle venu de l'extérieur, répond-elle aux besoins et aux réalités des sociétés africaines, ou doit-elle être adaptée pour y être efficace ?
\end{itemize}

\rubrique{2. Plan détaillé}

\begin{enumerate}
\item[I.] \textbf{La démocratie, un régime souhaitable pour l'Afrique.} 1. Elle garantit les droits et les libertés. 2. Elle permet l'alternance pacifique et limite les abus. 3. Elle n'est pas étrangère aux traditions africaines.
\item[II.] \textbf{Une démocratie qui peine à s'enraciner.} 1. Une démocratie souvent formelle (élections contestées, constitutions modifiées). 2. Des conditions socio-économiques et sécuritaires difficiles. 3. Un modèle perçu comme imposé de l'extérieur.
\item[III.] \textbf{Pour une démocratie adaptée aux réalités africaines.} 1. Intégrer les valeurs et institutions traditionnelles. 2. Une démocratie de proximité et de résultats. 3. Le rôle de l'éducation et de la société civile.
\end{enumerate}

\rubrique{3. Proposition de rédaction}

\partie{Introduction}

Depuis le début des années 1990, de nombreux pays africains ont adopté le multipartisme et les élections pluralistes, sous la pression des peuples et des partenaires extérieurs. Plus de trente ans plus tard, le bilan est contrasté : certains pays ont connu des alternances pacifiques, d'autres ont subi des élections contestées, des modifications constitutionnelles et, ces dernières années, un retour des coups d'État militaires, notamment au Sahel. D'où la question : la démocratie est-elle réellement la forme de gouvernement la plus adaptée aux pays africains ? N'est-elle qu'un modèle importé, inadapté aux réalités du continent, ou bien correspond-elle à une aspiration profonde des peuples africains ? Nous montrerons d'abord que la démocratie est un régime souhaitable pour l'Afrique, puis que son enracinement se heurte à de réelles difficultés, avant de plaider pour une démocratie adaptée aux réalités africaines.

\partie{I. La démocratie, un régime souhaitable pour l'Afrique}

\souspartie{1. Elle garantit les droits et les libertés}
La démocratie protège les libertés fondamentales : liberté d'expression, de presse, de réunion, d'association, droit de vote. Elle reconnaît la dignité de chaque citoyen et permet aux minorités de se faire entendre. Les peuples africains, qui ont connu la colonisation puis, souvent, des régimes de parti unique, aspirent à ces libertés, comme l'ont montré les conférences nationales du début des années 1990 (Bénin, 1990) ou l'insurrection populaire burkinabè d'octobre 2014 contre la révision de l'article 37 de la Constitution.

\souspartie{2. Elle permet l'alternance pacifique et limite les abus}
En organisant la conquête et la transmission du pouvoir par les urnes, la démocratie évite que le pouvoir ne se règle par les armes. Elle soumet les gouvernants au contrôle du parlement, de la justice et de la presse. Le Botswana, le Cap-Vert, le Ghana ou le Sénégal, qui ont connu des alternances pacifiques, montrent que la démocratie est possible et qu'elle favorise la stabilité. Comme le disait Winston Churchill, « la démocratie est le pire des régimes, à l'exception de tous les autres ».

\souspartie{3. Elle n'est pas étrangère aux traditions africaines}
Les sociétés africaines connaissaient des formes de délibération collective : l'arbre à palabres, où les décisions étaient prises après de longues discussions ; les conseils d'anciens qui limitaient le pouvoir du chef ; la Charte du Manden (XIII\textsuperscript{e} siècle), souvent citée comme l'une des premières déclarations des droits. L'Union africaine a d'ailleurs adopté en 2007 la Charte africaine de la démocratie, des élections et de la gouvernance.

\transition{La démocratie apparaît donc souhaitable. Pourtant, sa mise en œuvre en Afrique se heurte à de nombreux obstacles.}

\partie{II. Une démocratie qui peine à s'enraciner}

\souspartie{1. Une démocratie souvent formelle}
Dans de nombreux pays, les élections sont entachées de fraudes ou de contestations, les constitutions sont modifiées pour permettre aux dirigeants de se maintenir au pouvoir, les institutions de contrôle manquent d'indépendance. La démocratie se réduit alors à un rituel électoral sans véritable alternance, ce qui alimente la défiance des citoyens.

\souspartie{2. Des conditions socio-économiques et sécuritaires difficiles}
La pauvreté, l'analphabétisme et la dépendance économique limitent la participation éclairée des citoyens : le vote peut être acheté, l'information manipulée. L'ethnicisation de la vie politique transforme parfois les élections en affrontements communautaires. Dans les pays confrontés au terrorisme, l'urgence sécuritaire a conduit des militaires à prendre le pouvoir (Mali, Burkina Faso, Niger), avec le soutien d'une partie de la population déçue par les régimes élus.

\souspartie{3. Un modèle perçu comme imposé de l'extérieur}
La démocratie a parfois été introduite sous la pression des bailleurs de fonds, comme après le discours de La Baule (1990), qui liait l'aide française à la démocratisation. Les institutions copiées sur celles des anciennes puissances coloniales ne tiennent pas toujours compte des réalités locales. Certains y voient une forme d'ingérence et revendiquent une souveraineté politique.

\transition{Ces difficultés ne prouvent pas que la démocratie soit inadaptée à l'Afrique ; elles montrent qu'elle doit être enracinée dans les réalités africaines.}

\partie{III. Pour une démocratie adaptée aux réalités africaines}

\souspartie{1. Intégrer les valeurs et institutions traditionnelles}
Une démocratie africaine peut s'appuyer sur les valeurs de dialogue, de consensus et de solidarité, et associer les autorités coutumières et religieuses à la prévention des conflits et à la cohésion sociale. Il s'agit de combiner la légitimité moderne (le vote) et la légitimité traditionnelle (la sagesse, le consensus).

\souspartie{2. Une démocratie de proximité et de résultats}
La décentralisation rapproche le pouvoir des citoyens et leur permet de participer à la gestion de leurs affaires locales. Mais la démocratie doit aussi produire des résultats concrets — sécurité, santé, éducation, emploi — pour garder la confiance des peuples. Une bonne gouvernance, transparente et responsable, est la meilleure défense de la démocratie.

\souspartie{3. Le rôle de l'éducation et de la société civile}
Enfin, une démocratie solide a besoin de citoyens informés et engagés. L'éducation civique, l'alphabétisation, des médias libres et responsables, des syndicats et des associations actives sont les gardiens de la démocratie.

\partie{Conclusion}

La démocratie reste la forme de gouvernement la plus respectueuse des droits et la plus capable d'assurer une alternance pacifique ; elle correspond à une aspiration profonde des peuples africains et n'est pas étrangère à leurs traditions. Son enracinement est cependant difficile : élections contestées, pauvreté, insécurité et sentiment d'un modèle imposé de l'extérieur fragilisent sa légitimité. La solution ne réside pas dans le rejet de la démocratie, mais dans son adaptation : une démocratie enracinée dans les valeurs africaines, proche des citoyens et efficace dans la réponse à leurs besoins. Le défi pour l'Afrique est d'inventer une démocratie qui lui ressemble.

\end{spacing}
\end{document}
